% -----------------------------------------------------------------------------
% Elementos pré-textuais: resumo, abstract, lista de siglas e lista de símbolos
% -----------------------------------------------------------------------------

% ---
% Resumo em língua portuguesa
% ---
\setlength{\absparsep}{18pt}
\begin{resumo}
Editais de fomento à pesquisa e à inovação são publicados de forma dispersa
nos portais de agências e instituições, em formatos heterogêneos e com prazos
curtos, o que dificulta acompanhar as oportunidades com inscrições abertas.
Este trabalho desenvolve e avalia um sistema que localiza, na Web aberta,
editais de fomento com inscrições em aberto. O sistema gera consultas
direcionadas a domínios institucionais, recupera os documentos candidatos por
meio de uma API de busca e os submete a uma cadeia de filtros heurísticos, que
descarta páginas de listagem, avalia a relevância temática e extrai o prazo de
inscrição, seguida de deduplicação e persistência em banco de dados. Uma
interface web apresenta os editais vigentes e o banco é atualizado
diariamente. Cada coleta é registrada em um \textit{snapshot} com data de
referência fixa, o que permite reprocessá-la com outras configurações e
reproduzir exatamente as decisões do sistema. Os parâmetros das heurísticas
foram calibrados em cinco etapas, sobre um conjunto rotulado de 47 itens:
correção da avaliação, análise de sensibilidade um-fator-por-vez, avaliação de
variantes estruturais, busca em grade com validação cruzada e
\textit{bootstrap}, e congelamento da configuração, com as decisões tomadas por
um critério lexicográfico que prioriza não aceitar itens errados e não perder
editais. A configuração final foi avaliada uma única vez sobre um conjunto
independente de 26 itens, coletado em outra data. A análise de sensibilidade
revelou três defeitos de implementação invisíveis nas métricas agregadas,
entre eles uma extração de prazo cuja pontuação de frases nunca atuava. No
conjunto de teste, os editais descartados indevidamente caíram de 5 para 1
entre 14, com intervalo de confiança de 95\% da diferença de $-8$ a $-1$, e a
extração de prazo passou de nenhuma para 4 das 12 datas corretas, sem aumento
dos itens errados entregues ao usuário. A precisão do resultado final não
melhorou, pois o sistema ainda aceita páginas de listagem, o que indica a
necessidade de uma etapa que identifique se a página é um edital. Além do
sistema, o trabalho contribui com um procedimento de calibração reprodutível
e com o uso da análise de sensibilidade como instrumento de diagnóstico de
sistemas baseados em regras.

\vspace{\onelineskip}

\noindent
\textbf{Palavras-chave}: editais de fomento; recuperação de informação;
heurísticas; calibração de parâmetros; análise de sensibilidade.
\end{resumo}

% ---
% Resumo em língua estrangeira
% ---
\begin{resumo}[Abstract]
\begin{otherlanguage*}{english}
Calls for research and innovation funding are published in a scattered way
across the portals of agencies and institutions, in heterogeneous formats and
with short deadlines, which makes it hard to keep track of the opportunities
open for application. This work develops and evaluates a system that locates,
on the open Web, funding calls whose applications are open. The system
generates queries targeted at institutional domains, retrieves candidate
documents through a search API, and submits them to a chain of heuristic
filters, which discards listing pages, assesses topical relevance and extracts
the application deadline, followed by deduplication and storage in a
database. A web interface presents the open calls, and the database is updated
daily. Each collection is stored as a snapshot with a fixed reference date,
which allows it to be reprocessed with other configurations while reproducing
the system's decisions exactly. The heuristic parameters were calibrated in
five stages on a labeled set of 47 items: correction of the evaluation,
one-at-a-time sensitivity analysis, evaluation of structural variants, grid
search with cross-validation and bootstrap, and freezing of the configuration,
with decisions made by a lexicographic criterion that gives priority to not
accepting wrong items and not losing calls. The final configuration was
evaluated only once on an independent set of 26 items, collected on a
different date. The sensitivity analysis revealed three implementation defects
that were invisible in the aggregate metrics, among them a deadline extraction
whose sentence scoring never took effect. On the test set, the calls wrongly
discarded fell from 5 to 1 out of 14, with a 95\% confidence interval for the
difference of $-8$ to $-1$, and deadline extraction went from none to 4 of the
12 correct dates, with no increase in the wrong items delivered to the user.
The precision of the final result did not improve, since the system still
accepts listing pages, which points to the need for a step that decides
whether a page is a call. Besides the system, the work contributes a
reproducible calibration procedure and the use of sensitivity analysis as a
diagnostic tool for rule-based systems.

\vspace{\onelineskip}

\noindent
\textbf{Keywords}: funding calls; information retrieval; heuristics;
parameter calibration; sensitivity analysis.
\end{otherlanguage*}
\end{resumo}

% ---
% Lista de abreviaturas e siglas
% ---
\begin{siglas}
  \item[API]   Interface de Programação de Aplicações (\textit{Application
               Programming Interface})
  \item[CSS]   Folhas de Estilo em Cascata (\textit{Cascading Style Sheets})
  \item[CSV]   Valores Separados por Vírgula (\textit{Comma-Separated Values})
  \item[DSR]   \textit{Design Science Research}
  \item[FN]    Falso Negativo
  \item[FP]    Falso Positivo
  \item[HTML]  Linguagem de Marcação de Hipertexto (\textit{HyperText Markup
               Language})
  \item[HTTP]  Protocolo de Transferência de Hipertexto (\textit{HyperText
               Transfer Protocol})
  \item[HTTPS] Protocolo de Transferência de Hipertexto Seguro
               (\textit{HyperText Transfer Protocol Secure})
  \item[IC]    Intervalo de Confiança
  \item[JSON]  Notação de Objetos JavaScript (\textit{JavaScript Object
               Notation})
  \item[MD5]   Algoritmo de resumo de mensagem 5 (\textit{Message-Digest
               Algorithm 5})
  \item[OAT]   Um fator por vez (\textit{One-At-a-Time})
  \item[ORM]   Mapeamento Objeto-Relacional (\textit{Object-Relational
               Mapping})
  \item[PDF]   Formato de Documento Portátil (\textit{Portable Document Format})
  \item[RF]    Requisito Funcional
  \item[RNF]   Requisito Não Funcional
  \item[SHA]   Algoritmo de Resumo Seguro (\textit{Secure Hash Algorithm})
  \item[SVM]   Máquina de Vetores de Suporte (\textit{Support Vector Machine})
  \item[TCC]   Trabalho de Conclusão de Curso
  \item[URL]   Localizador Uniforme de Recursos (\textit{Uniform Resource
               Locator})
  \item[UTC]   Tempo Universal Coordenado (\textit{Coordinated Universal Time})
  \item[UTF-8] Formato de Transformação Unicode de 8 bits (\textit{8-bit
               Unicode Transformation Format})
  \item[VN]    Verdadeiro Negativo
  \item[VP]    Verdadeiro Positivo
\end{siglas}

% ---
% Lista de símbolos (na ordem de aparecimento no texto)
% ---
\begin{simbolos}
  \item[$P$]                      Precisão
  \item[$R$]                      Revocação
  \item[$f_1, \dots, f_n$]        Critérios de um ordenamento lexicográfico
  \item[$a$, $b$]                 Alternativas comparadas pelo critério
                                  lexicográfico
  \item[$L$]                      Pontuação de listagem
  \item[$\theta_L$]               Limiar de classificação como listagem
  \item[$\mathbf{1}[\cdot]$]      Função indicadora: vale 1 se a condição é
                                  verdadeira e 0 caso contrário
  \item[$w_k$]                    Peso da palavra de listagem na URL
  \item[$w_p$]                    Peso do padrão de listagem na URL
  \item[$w_d$]                    Peso de cada faixa de contagem de datas
  \item[$n_{\mathrm{datas}}$]     Número de datas no formato DD/MM/AAAA no texto
  \item[$d_{\mathrm{low}}$]       Primeira faixa de contagem de datas
  \item[$d_{\mathrm{high}}$]      Segunda faixa de contagem de datas
  \item[$w_l$]                    Peso da baixa ocorrência de termos de edital
  \item[$n_{\mathrm{det}}$]       Número de ocorrências de termos típicos de
                                  edital individual
  \item[$m$]                      Mínimo de ocorrências de termos de edital
  \item[$w_e$]                    Peso do plural ``editais''
  \item[$w_c$]                    Peso do plural ``chamadas públicas''
  \item[$K^{+}$, $K^{-}$]         Vocabulários positivo e negativo de relevância
  \item[$V$]                      Pontuação de relevância
  \item[$k$]                      Termo do vocabulário de relevância
  \item[$v_k$]                    Peso do termo $k$ do vocabulário de relevância
  \item[$n_k$]                    Número de ocorrências do termo $k$
  \item[$N$]                      Número total de palavras do texto
  \item[$\rho$]                   Densidade de termos positivos no texto
  \item[$B(\rho)$]                Bônus de densidade da relevância
  \item[$S$]                      Pontuação de uma frase candidata a prazo
  \item[$n_{\mathrm{pos}}$]       Número de termos indicativos de prazo
                                  distintos na frase
  \item[$n_{\mathrm{neg}}$]       Número de termos negativos distintos na frase
  \item[$F_{\beta}$]              Medida F com peso $\beta$
  \item[$\beta$]                  Peso relativo da revocação em relação à
                                  precisão na medida F
  \item[$F_{0{,}5}$, $F_1$, $F_2$] Medidas F com $\beta$ igual a 0,5, 1 e 2
  \item[$n$]                      Número de itens rotulados
  \item[$i$]                      Índice de um item rotulado
  \item[$p_i^{o}$, $p_i^{c}$]     Indicador de edital perdido do item $i$ nas
                                  configurações original e calibrada
  \item[$e_i^{o}$, $e_i^{c}$]     Indicador de item errado aceito do item $i$
                                  nas configurações original e calibrada
  \item[$\Delta_p$, $\Delta_e$]   Diferenças observadas de editais perdidos e
                                  de itens errados aceitos
  \item[$M$]                      Número de reamostras do \textit{bootstrap}
  \item[$j$]                      Posição de um sorteio em uma reamostra
  \item[$i_j^{*}$]                Índice do item sorteado na posição $j$ de
                                  uma reamostra
  \item[$\Delta_p^{*}$, $\Delta_e^{*}$] Diferenças calculadas em uma reamostra
\end{simbolos}
